\clearpage
\SourceLeaf{241}%
\begin{center}
{\Large\textsc{[Bolyai und Lobatschewsky.]}}



\emph{[\textup{\textsc{Gerling}} an \textup{\textsc{Gauss}}. Marburg, 18.--21. December 1843.]}
\end{center}


\{\mbox{.\ .\ .\ .\ .\ .\ .} Ein Paar Ungarsche Stipendiaten aus Maros Vásárhely brachten ein
Präsentations-Schreiben von dortiger Universität, was mir auffiel, weil ich unter
den übrigen Unterschriften auch den Namen \textsc{Wolfgangus Bolyai} fand. Sie
übergaben mir nun in diesen Tagen auch einen Privat-Brief von einem andern dortigen
Professor, der früher mein Zuhörer gewesen, und diese Gelegenheit benutzte ich, um
zur Mittheilung an Sie mich nach W. B. zu erkundigen. Sie schildern ihn als einen
ziemlich wohlbeleibten rüstigen und freundlichen Greis, der in der Mitte seiner
Gärten ein heiteres Alter zu verleben scheine. Seine Wirksamkeit als akademischer
Lehrer sei ununterbrochen, aber nur dadurch gehemmt, dass es den meisten Zuhörern zu
schwer falle ihm zu folgen. Sein Sohn (von welchem Sie mir mittheilten, dass er ein
gutes Buch geschrieben) sei als erst 32jähriger Officier als Rittmeister pensionirt
worden, und lebe, wenn ich recht verstanden habe, bei seinem Vater. --- Muthmasslich
hat er mit mathematischen Ketzereien also Anstoss erregt. --- Jener ehemalige
Zuhörer hatte mir damals versprochen, Anstalt zu machen, dass ich die Bücher von
Vater und Sohn (denn auch vom Vater habe ich einmal ein Buch bei Ihnen gesehen, was
Sie interessant nannten) hierher bekäme. Er scheint diess aber in Gnaden vergessen
zu haben, denn sein Brief enthält kein Wort davon. Auch kann ich nirgends (selbst in
Berlin nicht) die Titel mir ausforschen. --- Da nun die jetzt neu angekommenen
jungen Leute sich erbieten, für die Sache zu sorgen; so möchte ich Sie wohl bitten,
mir diese Titel gelegentlich aufschreiben zu lassen. \mbox{.\ .\ .\ .\ .\ .\ .}\}


